%% =========================================================
%% Chapter: Matrix Product Operator Algebras
%% Symmetry theory of matrix product operator algebras built on the
%% multi-block asymmetric compression theorem of Chapter ch:asymmetric_ft:
%% fusion algebras and their actions, fusion and action tensors, their gauge
%% freedom, associators and F-symbols, cohomological obstructions, and the
%% relation to weak Hopf algebras, dualities, and density operators.
%% Planned Lean modules:
%%   TNLean/MPS/Symmetry/MPOSymmetry/{Defs,NIMRep,Character,Examples}.lean
%%     (open PR #7963; the Lean names of Section sec:mpoalg_symmetry are kept
%%     in comments until that PR merges),
%%   TNLean/MPS/Symmetry/MPOSymmetry/FusionTensors.lean,
%%   TNLean/MPS/Symmetry/MPOSymmetry/Associator.lean.
%% =========================================================

\chapter{Matrix Product Operator Algebras}%
\label{ch:mpo_algebras}

A matrix product operator algebra is a finite family of matrix product
operators whose periodic operators multiply according to length-independent
fusion rules with nonnegative integer coefficients
\cite{Bultinck2017Anyons,GarreRubio2023MPOSym,Molnar2022MPOAlgebrasI}. This
chapter develops its symmetry theory along the lines of the on-site theory of
Chapter~\ref{ch:symmetry}: a symmetry is defined, it is transferred to the
virtual level, the virtual data are shown to be unique up to an explicit gauge
freedom, the gauge-invariant class is extracted, and a nontrivial class is
shown to obstruct symmetric states. The transfer to the virtual level rests on
the multi-block asymmetric compression theorem
(Theorem~\ref{thm:asym_multiblock}). The stacked tensor of a product of two
matrix product operators, or of an operator acting on a state, is never in
canonical form and in general admits no sitewise intertwiner with the tensors
of the fusion products (Remark~\ref{rem:asym_zipper}), so the Fundamental
Theorem for tensors in canonical form does not apply to it; the asymmetric
theorem, which requires equality of periodic traces only, does.

Throughout, $T_a$ denotes a matrix product operator tensor of bond dimension
$\chi_a$, $O_a^{(N)}$ its periodic operator on $N$ sites, and
$\ket{V^{(N)}(A)}$ the periodic vector of a matrix product state tensor $A$.
For two operator tensors, $T_aT_b$ denotes the stacked tensor
$(T_aT_b)^{ij}=\sum_mT_a^{im}\otimes T_b^{mj}$ of bond dimension
$\chi_a\chi_b$, and $T_a\cdot A$ denotes the tensor
$(T_a\cdot A)^i=\sum_jT_a^{ij}\otimes A^j$. A reduction of a tensor $B$ onto a
tensor $C$ is a pair of maps $(W,V)$ with $WV=\Id$ and $WB^wV=C^w$ for every
word $w$ (Definition~\ref{def:mps_rectangular_reduction}); we call $W$ the
compression and $V$ the embedding.

%% ---------------------------------------------------------
\section{Matrix product operator symmetries}%
\label{sec:mpoalg_symmetry}

% Source anchor: arXiv:2203.12563, Papers/2203.12563/REsubmission.tex,
%   lines 321--362 (algebras and fusion rules), 429--460 and 567--568
%   (symmetric states and the periodic-boundary action), 564--565 (the
%   multiplicity relation), 660 (unit and inverses).
% **Scope restriction (periodic boundary):** the source defines symmetry by
%   invariance of the arbitrary-boundary subspace (lines 431--434); this
%   section uses only its periodic-boundary consequence (lines 567--568).
%   Documented in docs/paper-gaps/glm23_mpo_symmetric_mps_scope.tex, added by
%   PR #7963; cite that note (bibliography key gap-colon plus its slug) once the
%   note is on main.
An invertible on-site symmetry acts on a matrix product state through a group
representation on the physical index. A matrix product operator symmetry is
given instead by a finite family of operators whose periodic operators
multiply according to fusion rules~\cite{GarreRubio2023MPOSym}. The source
defines the symmetry by invariance of the states with arbitrary boundary
conditions; the statements below use only its consequence for periodic
boundary conditions~\cite[Section~3]{GarreRubio2023MPOSym}.

% The planned Lean names in this section belong to the open PR #7963 and are
% kept in comments until it merges; leanok tags are added after checking each
% statement against the merged code.

\begin{definition}[Matrix product operator fusion algebra]%
	\label{def:mpo_fusion_algebra}
	% Planned Lean: MPOTensor.IsMPOFusionAlgebra
	\uses{def:mpo_tensor, def:mpo_family}
	A finite family $(T_a)_{a\in I}$ of matrix product operator tensors with
	structure constants $N_{ab}^c\in\N$ is a \emph{fusion algebra} if, for
	all $a,b\in I$ and $N\geq1$,
	\begin{align}
		O_a^{(N)}O_b^{(N)}
		 & =\sum_{c\in I}N_{ab}^c\,O_c^{(N)}.
		\label{eq:mpo_fusion_algebra}
	\end{align}
	This is the fusion rule of \cite[Section~2]{GarreRubio2023MPOSym} and
	\cite{Bultinck2017Anyons}.
	% Source anchor: arXiv:2203.12563, lines 361--362.
\end{definition}

\begin{definition}[Unit and invertible labels]%
	\label{def:fusion_unit_invertible}
	% Planned Lean: MPOTensor.IsFusionUnit, MPOTensor.IsInvertibleLabel
	\uses{def:mpo_fusion_algebra}
	A label $e$ is a \emph{unit} of the structure constants if
	$N_{eb}^c=N_{be}^c=\delta_{bc}$ for all $b,c$. A label $a$ is
	\emph{invertible} with respect to $e$ if there is a label $b$ with
	$N_{ab}^c=N_{ba}^c=\delta_{ce}$ for all $c$. The operator $O_e^{(N)}$ need
	not be the identity; outside the on-site case it is a projector
	\cite[Section~4]{GarreRubio2023MPOSym}.
	% Source anchor: arXiv:2203.12563, line 660.
\end{definition}

\begin{definition}[Symmetric family of matrix product states]%
	\label{def:mpo_symmetric_family}
	% Planned Lean: MPOTensor.IsMPOSymmetricFamily, MPOTensor.IsMPOSymmetric
	\uses{def:mpo_fusion_algebra, def:mps_tensor}
	A finite family $(A_x)_{x\in K}$ of matrix product state tensors is
	\emph{symmetric} under $(T_a)_{a\in I}$ with coefficients
	$M_{a,x}^y\in\C$ if, for all $a$, $x$ and $N\geq1$,
	\begin{align}
		O_a^{(N)}\ket{V^{(N)}(A_x)}
		 & =\sum_{y\in K}M_{a,x}^y\ket{V^{(N)}(A_y)}.
		\label{eq:mpo_symmetric_family}
	\end{align}
	A single tensor $A$ is symmetric with eigenvalues $c_a$ if
	$O_a^{(N)}\ket{V^{(N)}(A)}=c_a\ket{V^{(N)}(A)}$ for all $a$ and $N\geq1$.
	The coefficients do not depend on $N$.
	% Source anchor: arXiv:2203.12563, lines 567--568 (restated at line 1801).
\end{definition}

\begin{definition}[Fusion character and nonnegative integer representation]%
	\label{def:fusion_character_nimrep}
	% Planned Lean: MPOTensor.IsFusionCharacter, MPOTensor.IsNIMRep
	\uses{def:fusion_unit_invertible}
	A \emph{fusion character} with unit $e$ is a function $m$ on labels with
	$m_e=1$ and $m_am_b=\sum_cN_{ab}^cm_c$. A \emph{nonnegative integer
		representation} on a finite set $K$ is a family of matrices
	$M_{a,x}^y\in\N$ with
	\begin{align}
		\sum_cN_{ab}^cM_{c,x}^y
		 & =\sum_{z\in K}M_{b,x}^zM_{a,z}^y.
		\label{eq:mpo_nimrep}
	\end{align}
	% Source anchor: arXiv:2203.12563, lines 564--565 (the relation for the
	%   multiplicities); the fusion character is not used by the source.
\end{definition}

\begin{theorem}[The action coefficients form a nonnegative integer representation]%
	\label{thm:mpo_symmetric_nimrep}
	% Planned Lean: MPOTensor.exists_nat_eq_of_isMPOSymmetricFamily,
	%   MPOTensor.sum_fusion_mul_eq_sum_mul_of_isMPOSymmetricFamily,
	%   MPOTensor.exists_isNIMRep_of_isMPOSymmetricFamily
	\uses{def:mpo_symmetric_family, def:fusion_character_nimrep, def:normal}
	Let $(A_x)_{x\in K}$ be normal tensors of positive bond dimension whose
	periodic vectors are linearly independent at some length $N_0\geq1$,
	symmetric under a fusion algebra with coefficients $M_{a,x}^y$. Then
	every $M_{a,x}^y$ is a nonnegative integer and (\ref{eq:mpo_nimrep})
	holds. This is the integrality and associativity of the action
	multiplicities in \cite[Section~3]{GarreRubio2023MPOSym}.
	% Source anchor: arXiv:2203.12563, line 460 and lines 564--565.
\end{theorem}

\begin{proof}
	\uses{thm:asymex_fib_anomaly_trace_integral}
	Fix $a$ and $x$. The word traces of $T_a\cdot A_x$ are
	$\sum_yM_{a,x}^y\tr(A_y^w)$ on nonempty words $w$. The word modules of
	the $A_y$ are simple and pairwise non-isomorphic, since isomorphic word
	modules have equal periodic vectors. An element of the augmentation ideal
	of the word algebra acts on the block $y_0$ as a matrix unit and
	annihilates the other blocks; its powers have constant trace
	$M_{a,x}^{y_0}$ on $T_a\cdot A_x$, so $M_{a,x}^{y_0}$ is a nonnegative
	integer (Theorem~\ref{thm:asymex_fib_anomaly_trace_integral}). Applying
	(\ref{eq:mpo_fusion_algebra}) and (\ref{eq:mpo_symmetric_family}) twice
	to $O_a^{(N_0)}O_b^{(N_0)}\ket{V^{(N_0)}(A_x)}$ and comparing
	coefficients of the independent vectors gives (\ref{eq:mpo_nimrep}).
\end{proof}

\begin{theorem}[A single symmetric normal tensor gives a fusion character]%
	\label{thm:mpo_symmetric_character}
	% Planned Lean: MPOTensor.exists_isFusionCharacter_of_isMPOSymmetric
	\uses{def:mpo_symmetric_family, def:fusion_character_nimrep}
	If a normal tensor of positive bond dimension is symmetric under a fusion
	algebra with eigenvalues $c_a$ and $c_e=1$, then $c_a\in\N$ for every $a$
	and $a\mapsto c_a$ is a fusion character.
\end{theorem}

\begin{proof}
	\uses{thm:asymex_fib_anomaly_integral_action}
	Each $c_a$ is a nonnegative integer by
	Theorem~\ref{thm:asymex_fib_anomaly_integral_action}. Applying
	(\ref{eq:mpo_fusion_algebra}) to the nonzero periodic vector at an
	injective length gives $c_ac_b=\sum_cN_{ab}^cc_c$.
\end{proof}

\begin{corollary}[Rank-one obstruction]%
	\label{cor:mpo_symmetric_no_go}
	% Planned Lean: MPOTensor.not_isMPOSymmetric_of_forall_not_isFusionCharacter
	\uses{thm:mpo_symmetric_character}
	If the structure constants admit no fusion character with values in
	$\N$, no normal tensor of positive bond dimension is symmetric with
	$c_e=1$. This obstruction sees only the fusion ring; the cohomological
	obstruction of Theorem~\ref{thm:mpoalg_no_symmetric_normal}, which also
	excludes group-like algebras with a nontrivial three-cocycle, is a
	different statement.
\end{corollary}

\begin{proof}
	\uses{thm:mpo_symmetric_character}
	Theorem~\ref{thm:mpo_symmetric_character} produces a fusion character
	with values in $\N$ from any such tensor.
\end{proof}

\begin{theorem}[Invertible labels]%
	\label{thm:mpo_invertible_labels}
	% Planned Lean: MPOTensor.IsFusionCharacter.eq_one_of_isInvertibleLabel,
	%   MPOTensor.mpo_mulVec_eq_of_isInvertibleLabel,
	%   MPOTensor.IsNIMRep.exists_eq_one_of_isInvertibleLabel,
	%   MPOTensor.IsNIMRep.exists_equiv_of_isInvertibleLabel,
	%   MPOTensor.exists_equiv_mpo_mulVec_eq_of_isInvertibleLabel
	\uses{thm:mpo_symmetric_nimrep, thm:mpo_symmetric_character,
		def:fusion_unit_invertible}
	A fusion character with values in $\N$ is $1$ on every invertible label,
	so an invertible label fixes the periodic vectors of a single symmetric
	normal tensor with $c_e=1$. For a family as in
	Theorem~\ref{thm:mpo_symmetric_nimrep} on which $O_e^{(N)}$ fixes every
	periodic vector, each invertible label $g$ permutes the blocks: there is a
	bijection $x\mapsto g\cdot x$ of the block labels with
	$O_g^{(N)}\ket{V^{(N)}(A_x)}=\ket{V^{(N)}(A_{g\cdot x})}$ for all
	$N\geq1$. This is the permutation action of
	\cite[Section~4]{GarreRubio2023MPOSym}.
	% Source anchor: arXiv:2203.12563, line 683.
\end{theorem}

\begin{proof}
	\uses{thm:mpo_symmetric_nimrep, thm:mpo_symmetric_character}
	From $m_am_b=m_e=1$ in $\N$ one gets $m_a=1$. For the family, the
	multiplicity matrices satisfy $M_aM_b=M_bM_a=M_e=\Id$ with nonnegative
	integer entries, so each row of $M_a$ is a standard basis vector, and the
	two relations make the resulting map of blocks a bijection.
\end{proof}

\begin{example}[Fibonacci]%
	\label{ex:mpo_symmetry_fibonacci}
	% Planned Lean: FibonacciCompression.isMPOFusionAlgebra_fibBlock,
	%   FibonacciCompression.isFusionUnit_fibFusion,
	%   FibonacciCompression.isMPOSymmetricFamily_fibNimTargets,
	%   FibonacciCompression.isNIMRep_fibFusion,
	%   FibonacciCompression.not_isFusionCharacter_fibFusion,
	%   FibonacciCompression.not_exists_normal_fibonacci_symmetric
	\uses{cor:mpo_symmetric_no_go, thm:asymex_fib_anomaly_fusion_algebra,
		thm:asymex_fib_anomaly_nim_rep}
	The Fibonacci operators of Section~\ref{sec:asymex_fibonacci} form a
	fusion algebra with unit $1$ and $\tau\times\tau=1+\tau$; the two normal
	states with $\tau\cdot x_1=x_\tau$ and $\tau\cdot x_\tau=x_1+x_\tau$ form
	a symmetric family whose action is the regular representation
	\cite[Section~6]{GarreRubio2023MPOSym}. Since $m^2=1+m$ has no solution
	in $\N$, no single normal tensor $A$ satisfies
	$O_1^{(N)}\ket{V^{(N)}(A)}=\ket{V^{(N)}(A)}$ and
	$O_\tau^{(N)}\ket{V^{(N)}(A)}=c\ket{V^{(N)}(A)}$ for all $N\geq1$.
	% Source anchor: arXiv:2203.12563, line 1993; arXiv:1511.08090, App. D.1.
\end{example}

\begin{example}[Anomalous cyclic group of order three]%
	\label{ex:mpo_symmetry_z3}
	% Planned Lean: Z3Anomalous.isMPOFusionAlgebra_z3Block,
	%   Z3Anomalous.isInvertibleLabel_z3Fusion
	\uses{def:fusion_unit_invertible}
	The operators $1,U,U^\dagger$ of
	Section~\ref{sec:asymex_z3anomalous} form a group-like fusion algebra in
	which every label is invertible. The fusion ring does not see the
	nontrivial three-cocycle of this representation; the obstruction to
	invariant normal states comes from
	Theorem~\ref{thm:mpoalg_no_symmetric_normal}.
\end{example}

%% ---------------------------------------------------------
\section{Fusion tensors and action tensors}%
\label{sec:mpoalg_virtual}

The fusion rule (\ref{eq:mpo_fusion_algebra}) is a statement about periodic
operators. Its local form is a decomposition of the stacked tensor $T_aT_b$
into copies of the blocks $T_c$, and that decomposition exists because the
two sides of (\ref{eq:mpo_fusion_algebra}) have equal periodic traces on the
pair alphabet.

\begin{definition}[Fusion tensors]%
	\label{def:mpoalg_fusion_tensors}
	\notready
	% Planned Lean: TNLean/MPS/Symmetry/MPOSymmetry/FusionTensors.lean,
	%   MPOTensor.FusionTensors (structure).
	\uses{def:mpo_fusion_algebra, def:mps_rectangular_reduction}
	Let $(T_a)_{a\in I}$ be a fusion algebra of normal tensors of positive
	bond dimension. A family of \emph{fusion tensors} for the pair $(a,b)$
	consists of maps
	\begin{align}
		W_{ab}^{c,\mu} & \colon\C^{\chi_a}\otimes\C^{\chi_b}\to\C^{\chi_c}, &
		V_{ab}^{c,\mu} & \colon\C^{\chi_c}\to\C^{\chi_a}\otimes\C^{\chi_b},
		\label{eq:mpoalg_fusion_maps}
	\end{align}
	for $c\in I$ and $1\leq\mu\leq N_{ab}^c$, such that the pairs are mutually
	biorthogonal, $W_{ab}^{c,\mu}V_{ab}^{d,\nu}=\delta_{cd}\delta_{\mu\nu}
		\Id_{\chi_c}$, each pair is a reduction of $T_aT_b$ onto $T_c$,
	\begin{align}
		W_{ab}^{c,\mu}(T_aT_b)^{\mathbf w}V_{ab}^{c,\mu}
		 & =T_c^{\mathbf w}
		\label{eq:mpoalg_fusion_reduction}
	\end{align}
	for every word $\mathbf w$ in the pair alphabet, and the remainder
	$(T_aT_b)^{ij}-\sum_{c,\mu}V_{ab}^{c,\mu}T_c^{ij}W_{ab}^{c,\mu}$ generates
	a nilpotent algebra. This is the decomposition of
	\cite[Section~2]{GarreRubio2023MPOSym} and
	\cite{Bultinck2017Anyons}, with the nilpotent remainder of
	\cite[Section~4.4]{GarreRubio2023MPOSym} in place of an exact identity.
	% Source anchor: arXiv:2203.12563, lines 361--391 (fusion tensors and
	%   orthogonality) and lines 1026--1030 (nilpotent off-diagonal blocks);
	%   arXiv:1511.08090, AnyonsPEPS.tex lines 237--251.
\end{definition}

\begin{theorem}[Existence of fusion tensors]%
	\label{thm:mpoalg_fusion_tensors_exist}
	\notready
	% Planned Lean: TNLean/MPS/Symmetry/MPOSymmetry/FusionTensors.lean,
	%   MPOTensor.IsMPOFusionAlgebra.exists_fusionTensors.
	% Existing Lean special cases:
	%   MPOTensor.exists_multiBlockCompression_mulTensor_of_mpo_eq_sum
	%   (Theorem thm:asym_mpo_product) and
	%   MPOTensor.GroupFamily.IsRepresentation.exists_fusionTensors
	%   (Corollary cor:asym_fusion_tensors_group).
	\uses{def:mpoalg_fusion_tensors, thm:asym_mpo_product}
	Every pair $(a,b)$ of labels of a fusion algebra of normal tensors of
	positive bond dimension admits a family of fusion tensors. This is the
	existence statement of \cite[Appendix~A]{GarreRubio2023MPOSym}, which the
	source proves under its closedness condition; here only the fusion rule
	(\ref{eq:mpo_fusion_algebra}) at periodic boundary conditions is
	assumed.
	% Source anchor: arXiv:2203.12563, line 361 and Appendix A (lines
	%   2305--2455).
\end{theorem}

\begin{proof}
	\uses{thm:asym_mpo_product}
	Apply Theorem~\ref{thm:asym_mpo_product} with the target family
	consisting of $N_{ab}^c$ copies of $T_c$ for each $c$, all with weight
	$1$. Its compression pairs, indexed by $(c,\mu)$, are mutually
	biorthogonal reductions with a nilpotent remainder.
\end{proof}

For a group-like algebra, with labels in a group $G$ and
$N_{gh}^k=\delta_{k,gh}$, the fusion tensors reduce to one pair
$(W_{g,h},V_{g,h})$ for each pair of group elements; for representations by
matrix product unitaries this is Corollary~\ref{cor:asym_fusion_tensors_group}.

\begin{definition}[Action tensors]%
	\label{def:mpoalg_action_tensors}
	\notready
	% Planned Lean: TNLean/MPS/Symmetry/MPOSymmetry/FusionTensors.lean,
	%   MPOTensor.ActionTensors (structure).
	\uses{def:mpo_symmetric_family, def:mps_rectangular_reduction}
	Let $(A_x)_{x\in K}$ be a symmetric family of normal tensors of positive
	bond dimension whose coefficients $M_{a,x}^y$ are nonnegative integers. A
	family of \emph{action tensors} for $(a,x)$ consists of maps
	$W_{ax}^{y,i}\colon\C^{\chi_a}\otimes\C^{D_x}\to\C^{D_y}$ and
	$V_{ax}^{y,i}\colon\C^{D_y}\to\C^{\chi_a}\otimes\C^{D_x}$, for $y\in K$
	and $1\leq i\leq M_{a,x}^y$, that are mutually biorthogonal reductions of
	$T_a\cdot A_x$ onto the blocks $A_y$ with a nilpotent remainder.
	% Source anchor: arXiv:2203.12563, lines 459--490 (action tensors and
	%   orthogonality) and lines 1062--1091 (periodic-boundary form).
\end{definition}

\begin{theorem}[Existence of action tensors]%
	\label{thm:mpoalg_action_tensors_exist}
	\notready
	% Planned Lean: TNLean/MPS/Symmetry/MPOSymmetry/FusionTensors.lean,
	%   MPOTensor.IsMPOSymmetricFamily.exists_actionTensors.
	% Existing Lean: MPOTensor.exists_isReduction_actTensor_of_isNormal,
	%   MPOTensor.exists_multiBlockCompression_actTensor_of_isNormal
	%   (Corollary cor:asym_action_tensors).
	\uses{def:mpoalg_action_tensors, cor:asym_action_tensors,
		thm:mpo_symmetric_nimrep}
	Under the hypotheses of Theorem~\ref{thm:mpo_symmetric_nimrep}, every
	pair $(a,x)$ admits a family of action tensors. This is the existence of
	action tensors in \cite[Section~3]{GarreRubio2023MPOSym}.
	% Source anchor: arXiv:2203.12563, lines 459--460 and 1065--1091.
\end{theorem}

\begin{proof}
	\uses{cor:asym_action_tensors, thm:mpo_symmetric_nimrep}
	By Theorem~\ref{thm:mpo_symmetric_nimrep} the coefficients are
	nonnegative integers, so (\ref{eq:mpo_symmetric_family}) writes
	$O_a^{(N)}\ket{V^{(N)}(A_x)}$ as a sum of periodic vectors of blocks
	repeated with multiplicities $M_{a,x}^y$. Corollary~\ref{cor:asym_action_tensors}
	applied to this family with unit weights gives the reductions.
\end{proof}

%% ---------------------------------------------------------
\section{Gauge freedom of fusion and action tensors}%
\label{sec:mpoalg_gauge}

\begin{theorem}[Uniqueness of fusion tensors up to multiplicity gauge]%
	\label{thm:mpoalg_fusion_gauge_uniqueness}
	\notready
	% Planned Lean: TNLean/MPS/Symmetry/MPOSymmetry/FusionTensors.lean,
	%   MPOTensor.FusionTensors.exists_multiplicity_gauge.
	% Multiplicity-free case: MPSTensor.IsReduction.exists_boundary_dressed_proportional
	%   (Theorem thm:mps_reduction_boundary_dressed_uniqueness).
	\uses{def:mpoalg_fusion_tensors,
		thm:mps_reduction_boundary_dressed_uniqueness}
	Let $(W_{ab}^{c,\mu},V_{ab}^{c,\mu})$ and
	$(\widetilde W_{ab}^{c,\mu},\widetilde V_{ab}^{c,\mu})$ be two families
	of fusion tensors for the same pair $(a,b)$, and let $N_0$ be a common
	nilpotency bound of their remainders. Then for each $c$ there is a unique
	invertible matrix $Y_{ab}^c\in\GL_{N_{ab}^c}(\C)$ such that, for every
	word $\mathbf w$ in the pair alphabet with $|\mathbf w|>2N_0$,
	\begin{align}
		\widetilde W_{ab}^{c,\mu}(T_aT_b)^{\mathbf w}
		 & =\sum_{\nu}(Y_{ab}^c)_{\mu\nu}\,W_{ab}^{c,\nu}(T_aT_b)^{\mathbf w},
		\notag                                                                                 \\
		(T_aT_b)^{\mathbf w}\widetilde V_{ab}^{c,\mu}
		 & =\sum_{\nu}\bigl((Y_{ab}^c)^{-1}\bigr)_{\nu\mu}\,
		                                                   (T_aT_b)^{\mathbf w}V_{ab}^{c,\nu}.
		\label{eq:mpoalg_fusion_gauge}
	\end{align}
	If $N_{ab}^c=1$, then $Y_{ab}^c$ is a nonzero scalar. This is the
	uniqueness up to gauge of \cite[Section~2]{GarreRubio2023MPOSym} and the
	basis change of \cite[Section~4]{Molnar2022MPOAlgebrasI}, in the
	boundary-dressed form of \cite[Theorem~22]{Molnar2018SemiInjective}.
	% Source anchor: arXiv:2203.12563, lines 416--424;
	%   arXiv:2204.05940, mpo.tex line 1725;
	%   arXiv:1706.07329, cornerproblem.tex lines 3156--3162.
\end{theorem}

\begin{proof}
	\uses{thm:mps_reduction_boundary_dressed_uniqueness}
	For $N_{ab}^c=1$ this is
	Theorem~\ref{thm:mps_reduction_boundary_dressed_uniqueness} applied to
	the two reductions of $T_aT_b$ onto $T_c$. In general, the argument of
	that theorem compares the cross products
	$\widetilde W_{ab}^{c,\mu}(T_aT_b)^{\mathbf w}V_{ab}^{c,\nu}$: normality
	of $T_c$ and the commutant of the word span force each of them to be a
	scalar $(Y_{ab}^c)_{\mu\nu}$ times $T_c^{\mathbf w}$, and the same
	comparison in the opposite direction produces the inverse matrix.
\end{proof}

\begin{theorem}[Uniqueness of action tensors up to multiplicity gauge]%
	\label{thm:mpoalg_action_gauge_uniqueness}
	\notready
	% Planned Lean: TNLean/MPS/Symmetry/MPOSymmetry/FusionTensors.lean,
	%   MPOTensor.ActionTensors.exists_multiplicity_gauge.
	\uses{def:mpoalg_action_tensors,
		thm:mps_reduction_boundary_dressed_uniqueness}
	Two families of action tensors for the same pair $(a,x)$ are related, in
	the boundary-dressed sense of (\ref{eq:mpoalg_fusion_gauge}), by a unique
	matrix $X_{ax}^y\in\GL_{M_{a,x}^y}(\C)$ for each $y$; if
	$M_{a,x}^y=1$ it is a nonzero scalar. This is the gauge freedom of the
	action tensors in \cite[Section~3.1]{GarreRubio2023MPOSym}.
	% Source anchor: arXiv:2203.12563, lines 595--602.
\end{theorem}

\begin{proof}
	\uses{thm:mps_reduction_boundary_dressed_uniqueness}
	The proof of Theorem~\ref{thm:mpoalg_fusion_gauge_uniqueness} applies
	verbatim to the tensor $T_a\cdot A_x$ and the blocks $A_y$.
\end{proof}

%% ---------------------------------------------------------
\section{The associator of a group-like algebra}%
\label{sec:mpoalg_associator}

Let $G$ be a group and $(T_g)_{g\in G}$ a fusion algebra of normal tensors of
positive bond dimension with $N_{gh}^k=\delta_{k,gh}$, so that
$O_g^{(N)}O_h^{(N)}=O_{gh}^{(N)}$ for every $N\geq1$. Choose fusion tensors
$(W_{g,h},V_{g,h})$ for every pair. The triple product $T_gT_hT_k$ has two
composite reductions onto $T_{ghk}$,
\begin{align}
	W^{(gh)k}
	 & =W_{gh,k}\,(W_{g,h}\otimes\Id_{\chi_k}),
	 &
	W^{g(hk)}
	 & =W_{g,hk}\,(\Id_{\chi_g}\otimes W_{h,k}),
	\label{eq:mpoalg_composite_compressions}
\end{align}
with the corresponding composite embeddings.

\begin{lemma}[Composite reductions]%
	\label{lem:mpoalg_composite_reduction}
	\notready
	% Planned Lean: TNLean/MPS/Symmetry/MPOSymmetry/Associator.lean,
	%   MPOTensor.isReduction_mulTensor_comp.
	\uses{def:mpoalg_fusion_tensors}
	If $(W,V)$ is a reduction of $T_aT_b$ onto $T_c$ and $(W',V')$ a
	reduction of $T_cT_d$ onto $T_e$, then
	$(W'(W\otimes\Id),(V\otimes\Id)V')$ is a reduction of $T_aT_bT_d$ onto
	$T_e$ with a nilpotent remainder; the same holds on the other side and for
	an operator tensor acting on a state tensor.
\end{lemma}

\begin{proof}
	For a word $\mathbf w=((i_1,j_1),\ldots,(i_N,j_N))$, the stacked product
	factors as
	$(W\otimes\Id)(T_aT_bT_d)^{\mathbf w}(V\otimes\Id)
		=\sum_{\mathbf m}W(T_aT_b)^{(\mathbf i,\mathbf m)}V\otimes
		T_d^{(\mathbf m,\mathbf j)}
		=(T_cT_d)^{\mathbf w}$, and applying $(W',V')$ gives $T_e^{\mathbf w}$.
	The remainder is block upper triangular with nilpotent diagonal blocks.
\end{proof}

\begin{definition}[Associator]%
	\label{def:mpoalg_associator}
	\notready
	% Planned Lean: TNLean/MPS/Symmetry/MPOSymmetry/Associator.lean,
	%   MPOTensor.GroupFusionAlgebra.associator.
	\uses{lem:mpoalg_composite_reduction,
		thm:mpoalg_fusion_gauge_uniqueness}
	The \emph{associator} of the chosen fusion tensors is the function
	$\omega\colon G^3\to\C^\times$ defined by
	\begin{align}
		W^{g(hk)}(T_gT_hT_k)^{\mathbf w}
		 & =\omega(g,h,k)\,W^{(gh)k}(T_gT_hT_k)^{\mathbf w}
		\label{eq:mpoalg_associator}
	\end{align}
	for all sufficiently long words $\mathbf w$. It exists and is nonzero by
	Theorem~\ref{thm:mpoalg_fusion_gauge_uniqueness} applied to the two
	composite reductions of Lemma~\ref{lem:mpoalg_composite_reduction}, and
	it is unique because $T_{ghk}^{\mathbf w}\neq0$ for some long word. This is
	the three-cocycle of \cite[Section~4.4]{GarreRubio2023MPOSym}.
	% Source anchor: arXiv:2203.12563, lines 1028--1060 (periodic-boundary
	%   definition, words of length at least ell) and lines 664--676.
\end{definition}

\begin{theorem}[The associator is a normalized three-cocycle]%
	\label{thm:mpoalg_associator_cocycle}
	\notready
	% Planned Lean: TNLean/MPS/Symmetry/MPOSymmetry/Associator.lean,
	%   MPOTensor.GroupFusionAlgebra.associator_isCocycle.
	\uses{def:mpoalg_associator, def:mpug_three_cocycle}
	The associator satisfies the three-cocycle relation
	(\ref{eq:mpug_three_cocycle}),
	\begin{align}
		\omega(gh,k,l)\,\omega(g,h,kl)
		 & =\omega(g,h,k)\,\omega(g,hk,l)\,\omega(h,k,l),
		\label{eq:mpoalg_pentagon_group}
	\end{align}
	which is the pentagon equation for a group-like algebra
	\cite[Section~4]{GarreRubio2023MPOSym}.
	% Source anchor: arXiv:2203.12563, lines 677--679.
\end{theorem}

\begin{proof}
	\uses{lem:mpoalg_composite_reduction}
	The quadruple product $T_gT_hT_kT_l$ has five composite reductions onto
	$T_{ghkl}$, one for each bracketing. Applying
	(\ref{eq:mpoalg_associator}) inside Lemma~\ref{lem:mpoalg_composite_reduction}
	relates adjacent bracketings by one value of $\omega$. The two paths from
	$((gh)k)l$ to $g(h(kl))$ give the two sides of
	(\ref{eq:mpoalg_pentagon_group}) times the same nonzero boundary-dressed
	word, which proves the identity.
\end{proof}

\begin{theorem}[Gauge dependence of the associator]%
	\label{thm:mpoalg_associator_gauge}
	\notready
	% Planned Lean: TNLean/MPS/Symmetry/MPOSymmetry/Associator.lean,
	%   MPOTensor.GroupFusionAlgebra.associator_cohomologousTo.
	\uses{def:mpoalg_associator, thm:mpoalg_fusion_gauge_uniqueness,
		def:mpug_fusion_gauge, def:mpug_cohomologous_to}
	Replacing the fusion tensors by $(\beta_{g,h}W_{g,h},
		\beta_{g,h}^{-1}V_{g,h})$ with $\beta\colon G^2\to\C^\times$ replaces
	$\omega$ by $\beta\mathbin{\triangleright}\omega$ of
	Definition~\ref{def:mpug_fusion_gauge}. Consequently the class of
	$\omega$ under the relation of
	Definition~\ref{def:mpug_cohomologous_to} does not depend on the choice
	of fusion tensors, and depends only on the family $(T_g)_{g\in G}$.
	This is the coboundary freedom of \cite[Section~4]{GarreRubio2023MPOSym}.
	% Source anchor: arXiv:2203.12563, lines 680--682 and 720--726.
\end{theorem}

\begin{proof}
	\uses{thm:mpoalg_fusion_gauge_uniqueness, thm:mpug_fusion_scalar_gauge}
	The rescaled pairs are again fusion tensors
	(Theorem~\ref{thm:mpug_fusion_scalar_gauge}). Substituting them into
	(\ref{eq:mpoalg_composite_compressions}) multiplies $W^{g(hk)}$ by
	$\beta_{g,hk}\beta_{h,k}$ and $W^{(gh)k}$ by $\beta_{g,h}\beta_{gh,k}$,
	which is (\ref{eq:mpug_coboundary}). By
	Theorem~\ref{thm:mpoalg_fusion_gauge_uniqueness} every other choice of
	fusion tensors is boundary-dressed equal to such a rescaling.
\end{proof}

\begin{definition}[F-symbols of a multiplicity-free algebra]%
	\label{def:mpoalg_f_symbols}
	\notready
	% Planned Lean: TNLean/MPS/Symmetry/MPOSymmetry/Associator.lean,
	%   MPOTensor.FusionTensors.fSymbol.
	\uses{def:mpoalg_fusion_tensors, lem:mpoalg_composite_reduction,
		thm:mpoalg_fusion_gauge_uniqueness}
	Let $(T_a)_{a\in I}$ be a fusion algebra of normal tensors with
	$N_{ab}^c\in\{0,1\}$ and fusion tensors $(W_{ab}^c,V_{ab}^c)$. For labels
	$a,b,c,d$, the composite compressions
	$W_{ec}^d(W_{ab}^e\otimes\Id)$ over admissible $e$ and
	$W_{af}^d(\Id\otimes W_{bc}^f)$ over admissible $f$ are two families of
	reductions of $T_aT_bT_c$ onto the same number of copies of $T_d$. The
	matrix $F^{abc}_d$ relating them,
	\begin{align}
		W_{af}^d(\Id\otimes W_{bc}^f)(T_aT_bT_c)^{\mathbf w}
		 & =\sum_e\,(F^{abc}_d)_{fe}\,W_{ec}^d(W_{ab}^e\otimes\Id)
			                                     (T_aT_bT_c)^{\mathbf w}
		\label{eq:mpoalg_f_move}
	\end{align}
	for all sufficiently long words $\mathbf w$, exists, is unique, and is
	invertible by Theorem~\ref{thm:mpoalg_fusion_gauge_uniqueness}. For a
	group-like algebra $F^{ghk}_{ghk}=\omega(g,h,k)$. This is the $F$-move of
	\cite{Bultinck2017Anyons} and
	\cite[Section~2]{GarreRubio2023MPOSym}, written for compressions and in
	boundary-dressed form.
	% Source anchor: arXiv:1511.08090, AnyonsPEPS.tex lines 247--251
	%   (F-move for embeddings); arXiv:2203.12563, lines 392--409.
\end{definition}

\begin{theorem}[Pentagon equation]%
	\label{thm:mpoalg_pentagon}
	\notready
	% Planned Lean: TNLean/MPS/Symmetry/MPOSymmetry/Associator.lean,
	%   MPOTensor.FusionTensors.fSymbol_pentagon.
	\uses{def:mpoalg_f_symbols}
	The $F$-symbols of a multiplicity-free fusion algebra of normal tensors
	satisfy the pentagon equation
	\begin{align}
		\sum_h(F^{abc}_g)_{fh}\,(F^{ahd}_e)_{gi}\,(F^{bcd}_i)_{hj}
		 & =(F^{fcd}_e)_{gj}\,(F^{abj}_e)_{fi}
		\label{eq:mpoalg_pentagon}
	\end{align}
	in the index convention of \cite{Bultinck2017Anyons}, and
	replacing the fusion tensors by $(\beta_{ab}^cW_{ab}^c,
		(\beta_{ab}^c)^{-1}V_{ab}^c)$ changes each $F^{abc}_d$ by the
	corresponding ratio of the $\beta$'s. For group-like algebras
	(\ref{eq:mpoalg_pentagon}) is (\ref{eq:mpoalg_pentagon_group}).
	% Source anchor: arXiv:1511.08090, AnyonsPEPS.tex lines 279--283;
	%   arXiv:2203.12563, lines 410--424; arXiv:2204.05940, mpo.tex line 1892.
	% The index placement in (eq:mpoalg_pentagon) must be rechecked against
	% the compression convention of (eq:mpoalg_f_move) when the Lean
	% statement is written.
\end{theorem}

\begin{proof}
	\uses{def:mpoalg_f_symbols, lem:mpoalg_composite_reduction}
	As for Theorem~\ref{thm:mpoalg_associator_cocycle}, compare the five
	bracketings of the reductions of $T_aT_bT_cT_d$ onto copies of $T_e$
	along the two paths of the pentagon, and use uniqueness of the matrix in
	(\ref{eq:mpoalg_f_move}).
\end{proof}

%% ---------------------------------------------------------
\section{Cohomological obstruction to symmetric states}%
\label{sec:mpoalg_obstruction}

\begin{definition}[L-symbols of an action]%
	\label{def:mpoalg_l_symbols}
	\notready
	% Planned Lean: TNLean/MPS/Symmetry/MPOSymmetry/Associator.lean,
	%   MPOTensor.GroupFusionAlgebra.lSymbol.
	\uses{def:mpoalg_associator, def:mpoalg_action_tensors,
		thm:mpo_invertible_labels, def:mpug_l_symbols}
	In the setting of Section~\ref{sec:mpoalg_associator}, let $(A_x)_{x\in
			K}$ be a symmetric family as in Theorem~\ref{thm:mpo_invertible_labels},
	so that $O_g^{(N)}\ket{V^{(N)}(A_x)}=\ket{V^{(N)}(A_{g\cdot x})}$, with
	action tensors $(W_{g,x},V_{g,x})$ reducing $T_g\cdot A_x$ onto
	$A_{g\cdot x}$. The tensor $T_gT_h\cdot A_x$ has two composite reductions
	onto $A_{gh\cdot x}$, and the \emph{$L$-symbol} $L^x_{g,h}\in\C^\times$ is
	defined by
	\begin{align}
		W_{gh,x}(W_{g,h}\otimes\Id)(T_gT_h\cdot A_x)^w
		 & =L^x_{g,h}\,W_{g,h\cdot x}(\Id\otimes W_{h,x})
		(T_gT_h\cdot A_x)^w
		\label{eq:mpoalg_l_symbol}
	\end{align}
	for all sufficiently long words $w$. This is the $L$-symbol of
	\cite[Section~4.4]{GarreRubio2023MPOSym}.
	% Source anchor: arXiv:2203.12563, lines 1091--1129 (periodic boundary),
	%   lines 700--716 (group case).
\end{definition}

\begin{theorem}[Compatibility of L-symbols with the associator]%
	\label{thm:mpoalg_l_compatible}
	\notready
	% Planned Lean: TNLean/MPS/Symmetry/MPOSymmetry/Associator.lean,
	%   MPOTensor.GroupFusionAlgebra.lSymbol_isCompatible.
	\uses{def:mpoalg_l_symbols, def:mpug_l_symbols, def:mpug_l_gauge}
	The $L$-symbols of Definition~\ref{def:mpoalg_l_symbols} are compatible
	with the associator in the sense of
	(\ref{eq:mpug_l_compatibility}),
	\begin{align}
		L^x_{g,hk}\,L^x_{h,k}
		 & =\omega(g,h,k)\,L^{k\cdot x}_{g,h}\,L^x_{gh,k},
		\label{eq:mpoalg_l_compatible}
	\end{align}
	and rescaling the fusion and action tensors by $\beta_{g,h}$ and
	$\gamma^x_g$ changes $(\omega,L)$ to
	$(\beta\mathbin{\triangleright}\omega,(\beta,\gamma)\mathbin{\triangleright}L)$
	of Definitions~\ref{def:mpug_fusion_gauge} and~\ref{def:mpug_l_gauge}.
	This is the coupled pentagon equation for a group-like algebra
	\cite[Section~4]{GarreRubio2023MPOSym}.
	% Source anchor: arXiv:2203.12563, lines 716--726 (pentagongroups and
	%   gdgroup).
\end{theorem}

\begin{proof}
	\uses{lem:mpoalg_composite_reduction, def:mpoalg_associator}
	The five composite reductions of $T_gT_hT_k\cdot A_x$ onto
	$A_{ghk\cdot x}$ corresponding to $g(h(k x))$, $g((hk)x)$, $(g(hk))x$,
	$((gh)k)x$ and $(gh)(kx)$ are related consecutively by $L^x_{h,k}$,
	$L^x_{g,hk}$, $\omega(g,h,k)$, $L^x_{gh,k}$ and $L^{k\cdot x}_{g,h}$;
	comparing the two paths gives (\ref{eq:mpoalg_l_compatible}). The gauge
	statement follows by substituting the rescaled tensors into
	(\ref{eq:mpoalg_associator}) and (\ref{eq:mpoalg_l_symbol}).
\end{proof}

\begin{theorem}[No symmetric normal tensor under a nontrivial associator]%
	\label{thm:mpoalg_no_symmetric_normal}
	\notready
	% Planned Lean: TNLean/MPS/Symmetry/MPOSymmetry/Associator.lean,
	%   MPOTensor.GroupFusionAlgebra.not_isMPOSymmetric_of_not_isTrivialGaugeClass;
	%   algebraic step TNLean.Algebra.LSymbol.IsCompatible.isTrivialGaugeClass_of_unique.
	\uses{thm:mpoalg_l_compatible, def:mpug_trivial_gauge_class,
		thm:mpoalg_associator_gauge}
	Let $(T_g)_{g\in G}$ be a group-like fusion algebra of normal tensors
	whose associator does not have trivial gauge class. Then there is no
	normal tensor $A$ of positive bond dimension with
	$O_g^{(N)}\ket{V^{(N)}(A)}=\ket{V^{(N)}(A)}$ for all $g\in G$ and
	$N\geq1$. This is the statement of
	\cite[Section~3.2]{GarreRubio2023MPOSym} that an injective state without
	multiplicity cannot be invariant under a nontrivial algebra, specialized
	to group-like algebras; for group-like algebras it goes back to
	\cite{ChenLiuWen2011CZX}.
	% Source anchor: arXiv:2203.12563, lines 607--636.
\end{theorem}

\begin{proof}
	\uses{thm:mpoalg_l_compatible, thm:mpoalg_action_tensors_exist}
	Action tensors exist by Theorem~\ref{thm:mpoalg_action_tensors_exist}
	with $K$ a single point fixed by every $g$. Then
	(\ref{eq:mpoalg_l_compatible}) reads
	$L_{g,hk}L_{h,k}=\omega(g,h,k)L_{g,h}L_{gh,k}$, that is,
	$\omega=dL$ in the notation of (\ref{eq:mpug_coboundary}). Hence
	$\omega=L\mathbin{\triangleright}1$ has trivial gauge class, a
	contradiction.
\end{proof}

%% ---------------------------------------------------------
\section{Weak Hopf algebras, dualities, and density operators}%
\label{sec:mpoalg_backbone}

\begin{remark}[Fusion tensors from a coproduct]%
	\label{rem:mpoalg_weak_hopf}
	For a matrix product operator representation of a cosemisimple weak Hopf
	algebra $\mathcal A$, the operators are indexed by irreducible
	representations $\psi_a$ of the dual algebra and the fusion tensors
	decompose the coproduct,
	$(\psi_a\otimes\psi_b)\circ\Delta(f)=\sum_{c,\mu}V_{ab}^{c\mu}\psi_c(f)
		W_{ab}^{c\mu}$ with $W_{ab}^{c\mu}V_{ab}^{d\nu}=
		\delta_{cd}\delta_{\mu\nu}\Id$ \cite[Section~4]{Molnar2022MPOAlgebrasI}.
	These are exact identities, and they are fusion tensors in the sense of
	Definition~\ref{def:mpoalg_fusion_tensors} with zero remainder;
	coassociativity of $\Delta$ is the associativity used in
	Definition~\ref{def:mpoalg_f_symbols}. Only this part of the dictionary
	between weak Hopf algebras and matrix product operators is in scope here.
	% Source anchor: arXiv:2204.05940, mpo.tex lines 1640--1660 (fusion
	%   tensors from the coproduct) and 1725--1892 (basis change, F-symbols,
	%   pentagon).
\end{remark}

\begin{remark}[Kramers--Wannier duality is an intertwiner]%
	\label{rem:mpoalg_kramers_wannier}
	The Kramers--Wannier operator of Section~\ref{sec:asymex_kw} does not
	belong to a fusion algebra on one lattice: by
	(\ref{eq:asymex_kw_fusion}) its square is $(\Id+\eta)T$ up to
	normalization, which contains the translation $T$ and is not a
	length-independent combination of the operators of a finite family.
	It is a matrix product operator intertwiner between two symmetric
	Hamiltonians, a bimodule between the symmetry algebras of the primal and
	dual chains, rather than a symmetry of one of them
	\cite{Lootens2023Dualities,Lootens2021MPOStringNet}. Fusion and action
	tensors for such intertwiners come from the same asymmetric compression
	theorem, applied to products of operators between different chains.
	% Source anchor: arXiv:2112.09091, _Hamiltonian.tex lines 22--28 and
	%   _Examples.tex lines 55--63.
\end{remark}

\begin{remark}[Density operators]%
	\label{rem:mpoalg_mpdo}
	Renormalization fixed points of matrix product density operators are
	described by matrix product operator algebras of $C^*$-weak Hopf algebras
	through the pulling-through identity of the canonical regular element
	\cite{RuizDeAlarcon2024MPOAlgebrasII}; the density-operator setting is
	that of Chapter~\ref{ch:mpdo}. The fusion tensors of
	Section~\ref{sec:mpoalg_virtual} apply there without change, since they
	use only equality of periodic traces.
	% Source anchor: arXiv:2204.06295, main.tex lines 461--470 (C*-weak Hopf
	%   algebra) and 802--818 (canonical regular element).
\end{remark}
